%%%%%%%%%%%%%%%%%%%%%%%%%%%%%%%%%%%%%%%%%
% University Assignment Title Page 
% LaTeX Template
% Version 1.0 (27/12/12)
%
% This template has been downloaded from:
% http://www.LaTeXTemplates.com
%
% Original author:
% WikiBooks (http://en.wikibooks.org/wiki/LaTeX/Title_Creation)
%
% License:
% CC BY-NC-SA 3.0 (http://creativecommons.org/licenses/by-nc-sa/3.0/)
% 
% Instructions for using this template:
% This title page is capable of being compiled as is. This is not useful for 
% including it in another document. To do this, you have two options: 
%
% 1) Copy/paste everything between \begin{document} and \end{document} 
% starting at \begin{titlepage} and paste this into another LaTeX file where you 
% want your title page.
% OR
% 2) Remove everything outside the \begin{titlepage} and \end{titlepage} and 
% move this file to the same directory as the LaTeX file you wish to add it to. 
% Then add \input{./title_page_1.tex} to your LaTeX file where you want your
% title page.
%
%%%%%%%%%%%%%%%%%%%%%%%%%%%%%%%%%%%%%%%%%
%\title{Title page with logo}
%----------------------------------------------------------------------------------------
%	PACKAGES AND OTHER DOCUMENT CONFIGURATIONS
%----------------------------------------------------------------------------------------

\documentclass[12pt]{article}
\usepackage[english]{babel}
\usepackage[utf8x]{inputenc}
\usepackage{amsmath}
\usepackage{graphicx}
\usepackage[colorinlistoftodos]{todonotes}

\begin{document}

\begin{titlepage}

\newcommand{\HRule}{\rule{\linewidth}{0.5mm}} % Defines a new command for the horizontal lines, change thickness here

\center % Center everything on the page


\includegraphics{logo.png}\\[.1cm] % Include a department/university logo - this will require the graphicx package

%----------------------------------------------------------------------------------------
%	HEADING SECTIONS
%----------------------------------------------------------------------------------------

\text{\LARGE Department of Computer Science and Engineering}\\[1.5cm] % Name of your university/college

%
\textsc{\Large Report on }\\[0.5cm] % Major heading such as course name
%----------------------------------------------------------------------------------------
%	TITLE SECTION
%----------------------------------------------------------------------------------------

\HRule \\[0.4cm]
{ \huge \bfseries ARA – The Ant-Colony Based Routing Algorithm for MANETs}\\[0.4cm] % Title of your document
{Mesut G¨unes, Udo Sorges, Imed Bouazizi}
\HRule \\[1cm]
 
%----------------------------------------------------------------------------------------
%	AUTHOR SECTION
%----------------------------------------------------------------------------------------

\begin{minipage}{0.5\textwidth}
\begin{flushleft} \large
\emph{Report Writers:}\\
Shantanu Sarkar\newline(0416052041) \newline % Your name
Mahmud Ahmed\newline(0416052027) \newline
Md. Mostafizur Rahman\newline(0416052032) 
\end{flushleft}
\end{minipage}
~
\begin{minipage}{0.4\textwidth}
\begin{flushright} \large
\emph{Supervisor:} \\
Dr. Ashikur  Rahman% Supervisor's Name
\end{flushright}
\end{minipage}\\[1cm]



% If you don't want a supervisor, uncomment the two lines below and remove the section above
%\Large \emph{Author:}\\
%John \textsc{Smith}\\[3cm] % Your name

%----------------------------------------------------------------------------------------
%	DATE SECTION
%----------------------------------------------------------------------------------------

{\large \today}\\[.5cm] % Date, change the \today to a set date if you want to be precise

%----------------------------------------------------------------------------------------
%	LOGO SECTION
%----------------------------------------------------------------------------------------


 
%----------------------------------------------------------------------------------------

\vfill % Fill the rest of the page with whitespace

\end{titlepage}

\tableofcontents

\newpage
\begin{abstract}

\end{abstract}

\section{Introduction}
Introduction of the report goes here. 


\section{Context and Problem Statement}


\section{Idea}

\section{Evaluation Metrics}
According to the Authors, the main feature of ARA is its low routing overhead and easy maintainability of routes between nodes in the topology. So, The evaluation metrics considered to measure the performance of ARA are:
\begin{itemize}
\item Delivery rate 
\item Routing overhead in terms of bits
\item Routing overhead in terms of packets
\end{itemize}

\section{Evaluation Process}
The performance of ARA was evaluated using simulation in terms of evaluation metrics mentioned earlier. 
The simulation was implemented in ns-2. Some important parameters of the simulation environment are:

\begin{itemize}
\item Simulation area  1500m×300m 
\item Maximum velocity of nodes 10 m/s
\item Simulation time  900 seconds
\item 10 Constant bit rate(CBR) connections
\item 7 different pause times\footnote{pause time indicates the mobility of the nodes} 0,30, 60, 120, 300, 600 and 900 seconds
\end{itemize}
Multiple simulations were run and the results were collected for evaluation. 


\section{Evaluation Results}

\section{Limitations of ARA}

\section{Future Works}

\section{Conclusion}

\end{document}